\documentclass{article}

\begin{document}
    \begin{center}
        \Large \textbf{Challenge 2: Shell} \\ [6pt]
        \normalsize
        Autor: Maximilian Jakob Maag B.Sc. \\
        Version: 1.0
    \end{center}
    
    \section{Intro}
    Egal ob Server oder Desktop, die Shell ist unter Linux das standard Instrument für
    alle Aufgaben. Es gibt verschiedene Shells, die bekannteste ist die Bash Shell. \\\\
    In der Shell können Befehle eingegeben werden, um Dateien zu erstellen, zu löschen, zu verschieben, um Software zu installieren, um Prozesse zu starten und zu stoppen, um Netzwerkverbindungen zu verwalten und vieles mehr. \\\\
    Es ist wichtig, die Shell zu verstehen, da sie ein mächtiges Werkzeug ist, um mit dem Betriebssystem zu interagieren.

    \section{Lernziele}
    Sie lernen die Shell zu bedienen. Nutzen Sie die Shell als Schweitzer Taschenmesser für die Interaktion jeder beliebigen Linux Distro.

    \section{Aufgaben}
    \subsection{Aufgabe 1}
    Öffnen Sie eine Shell/Terminal und geben Sie sich selbst root rechte innerhalb dieser Shell Session.

    \subsection{Aufgabe 2}
    Erstellen Sie eine Datei namens test.txt in ihrem Home Verzeichnis. \\
    Überprüfen Sie, ob die Datei erstellt wurde.

    \subsection{Aufgabe 3}
    Installieren Sie die beiden tools ls und less, falls nicht bereits vorhanden. \\
    Nutzen Sie eine Pipe und ein output redirect, um das Ergebniss von ls -l /etc in eine Textdatei in ihrem Home Verzeichnis zu schreiben.

    \subsection{Aufgabe 4}
    Öffnen Sie das Verzeichnis /var/log und geben Sie das das Ergebniss von PWD aus. \\
    Öffnen Sie nun ein anderes Verzeichnis und geben Sie PWD erneut aus. Was hat sich geändert und warum?

    \subsection{Aufgabe 5}
    Schreiben Sie ein simples Hello World Programm in C. Kompilieren Sie es und führen Sie es aus.
    Das Programm soll "Hello World" in der Shell ausgeben. \\
    Fügen Sie dieses Programm Ihrem Shell PATH hinzu. Erläutern Sie die Bedeutung dieser Änderung.

    \subsection{Aufgabe 6}
    Ändern Sie das C Programm aus Aufgabe 5 so, dass es die aktuelle Uhrzeit ausgibt. \\
    Fügen Sie eine Option zu ihrem Programm hinzu, wird es mit --24h aufgebrufen muss die Ausgabe z.B. 14:30 Uhr ausgeben. \\
    WIrd das Programm mit --12h aufgerufen soll die Ausgabe z.B. 2:30 PM sein. \\ \\
    Legen Sie zum schluss einen Alias an, damit das Programm mit dem Befehl "time" aufgerufen werden kann, mit der Option --12h.

    \subsection{Aufgabe 7}
    Nutzen Sie die manpage des Befehls "ls" um herauszufinden, was die Option -a bewirkt.
\end{document}